\section{Experimental Methodology}
\label{sec:methodology}
\label{sec:methodology-anchor}
The campaign is defined by a single JSON specification,
\texttt{experiments/five\_ue\_traffic\_scenarios\_v5.json}
(\texttt{protocol\_version}~8, SHA-256 \texttt{7f1d5c03\ldots fea025d}), and executed by
\texttt{experiments/experiment\_runner.py}. The specification fixes three traffic scenarios,
three arms, five seeds, the 106-PRB cell, the 100~ms control period, the 50~ms monitor
summary period, the 10~s rApp period, every phase duration, and every validity threshold.
Nothing in the protocol below is a per-run operator decision. At launch the runner copies the
specification over the static controller configuration --- the DDPG control period and the
rApp period in \texttt{period\_ms}, \texttt{monitor.summary\_period\_ms},
\texttt{operational\_min\_prb}/\texttt{operational\_max\_prb} into
\texttt{control.min\_prb}/\texttt{control.max\_prb}, the cell PRB count
\texttt{ddpg.cell\_prbs}, \texttt{ddpg.model\_version}, the channel-profile enable flag, and
the spec-derived \texttt{exploration\_noise\_decay\_steps} --- and writes the merged
configuration to \texttt{<run>/effective\_config.json}, which is the authoritative record of
what a run actually executed with.

\subsection{Compared control arms and their ablation}
Three arms are compared. Their code identifiers are \texttt{llm\_only}, \texttt{ddpg\_only},
and \texttt{llm\_guided\_ddpg}; in the text we call them LLM-only, DDPG-only (TD3-style), and
LLM-guided. The generated tables emitted by \texttt{scripts/build\_results.py} label the same
three arms ``LLM only'', ``TD3 only'', and ``LLM-guided TD3''.

\begin{itemize}
  \item \texttt{llm\_only} executes the rApp target. The published ratios are projected into
  the intersection of the A1 guidance bounds and the operational envelope, converted to an
  integer Slice-A budget, and commanded every 100~ms. The target itself is refreshed only
  every 10~s, so this arm holds one action between rApp updates. No actor--critic is used and
  no learner process is started.
  \item \texttt{ddpg\_only} executes the learned actor. It is an ablation of guidance, not of
  the task: it keeps the identical 24-dimensional state layout, the identical SLA-gated
  reward with the same weights, the identical intent semantics (priority slice, protected
  slice, floor $F$, and minimum priority ratio), and the identical calibration profile of its
  seed. What is removed is the LLM itself. In \texttt{state\_vector} the guidance flag is
  false, so both projected-target features are zeroed; the arm is bounded only by the
  operational envelope and never intersects the A1 \texttt{policy\_constraints}; the
  behavioral-cloning teacher term is disabled (\texttt{teacher\_mask}${}=0$, so the actor loss
  reduces to the deterministic policy-gradient term plus the saturation penalty); and the
  intent-triggered fallback to the LLM action does not exist.
  \item \texttt{llm\_guided\_ddpg} commands a convex blend
  $b_A=(1-w)\,b_A^{\mathrm{LLM}}+w\,b_A^{\mathrm{RL}}$ clipped to the active feasible
  interval. The learned weight $w$ starts at $0$, rises by $0.01$ every ten safety updates
  while the last 30 \emph{intent} outcomes are at least 95\% satisfied, is capped at
  \texttt{max\_ddpg\_weight}${}=0.4$, and is reduced by $0.1$ (and forced to $0$ for command
  purposes) after five consecutive \emph{intent} violations. The trigger is the
  intent-satisfaction flag and not SLA satisfaction: \texttt{update\_guided\_safety} is
  invoked with the per-transition intent-satisfied component, even though the fixed-length
  window it appends to is misleadingly named \texttt{sla\_history}. The distinction is
  material rather than cosmetic, because intent satisfaction is the strictly stronger of the
  two conditions defined below and the two diverge widely in the pilot
  (Section~\ref{sec:results}). During critic warm-up the commanded budget is additionally
  dithered by $\pm2$ PRBs around the teacher. The commanded action, not the raw actor output,
  is what enters replay.
\end{itemize}

Table~\ref{tab:proto-arms} summarizes the ablation. Because the objective and the state
layout are shared, a difference between \texttt{ddpg\_only} and \texttt{llm\_guided\_ddpg}
is \emph{not} attributable to a different task definition. It is not thereby attributable to
guidance alone: Section~\ref{sec:limitations} records a residual confound in episode
termination---the guided arm's terminal flag fires on every 10~s A1 policy version bump,
whereas the unguided arm's fires only at an intent switch, so the two arms bootstrap over
episodes of different length---and a single seed cannot separate policy effects from seed
effects.

\begin{table}[t]
\centering
\footnotesize
\setlength{\tabcolsep}{4pt}
\caption{Ablation across the three arms. ``Intent semantics'' covers the priority/protected
identities, the calibrated floor $F$, and the 55\% minimum priority ratio. The
intent-violation row is the fallback to the LLM action after five consecutive intent
violations; it is driven by intent satisfaction, not by SLA satisfaction.}
\label{tab:proto-arms}
\begin{tabular}{lccc}
\toprule
Component & LLM-only & DDPG-only & LLM-guided\\
\midrule
Intent semantics in reward      & yes & yes & yes\\
24-D state layout               & yes & yes & yes\\
Projected A1 target in state    & yes & zeroed & yes\\
A1 guidance bounds applied      & yes & no & yes\\
Operational 10--90\% envelope   & yes & yes & yes\\
Learned actor--critic           & no  & yes & yes\\
Teacher (BC) loss term          & --- & no  & yes\\
Intent-violation LLM fallback   & --- & no  & yes\\
Asynchronous learner started    & no  & yes & yes\\
\bottomrule
\end{tabular}
\end{table}

\subsection{Run protocol and step budget}
A run executes the phases of Table~\ref{tab:proto-phases}. Phase durations are converted to
control steps by \texttt{\_phase\_steps} as
$\operatorname{round}(\text{seconds}\times1000/\text{period})$; with the 100~ms period
this yields 1{,}000 calibration, 36{,}000 training, and $1{,}800+1{,}800$ evaluation steps,
i.e.\ 40{,}600 counted transitions per run. \texttt{\_validate\_run} compares the per-phase
transition counts in \texttt{experiment\_steps} against these exact values and aborts the run
on any mismatch, so a partially executed phase can never enter the results.

\begin{table}[t]
\centering
\footnotesize
\caption{Per-run phase budget at the 100~ms control period.}
\label{tab:proto-phases}
\begin{tabular}{lrrll}
\toprule
Phase & Dur.\ (s) & Steps & Learning & Active intent\\
\midrule
\texttt{settle}        & 30   & ---     & ---      & bootstrap\\
\texttt{calibration}   & 100  & 1{,}000 & replayed & bootstrap\\
\texttt{intent\_switch}& var. & var.    & training only & regenerating\\
\texttt{training}      & 3600 & 36{,}000& on       & alternating\\
\texttt{intent1\_eval} & 180  & 1{,}800 & frozen   & Intent~1\\
\texttt{intent2\_eval} & 180  & 1{,}800 & frozen   & Intent~2\\
\midrule
counted total          & 4060 & 40{,}600 & & \\
\bottomrule
\end{tabular}
\end{table}

Before settling, the runner installs a bootstrap intent (Intent~1 with a 10~Mbps placeholder
floor and no policy constraints), waits for the first rApp guidance, verifies UE reachability,
and starts the iperf3 servers and clients. It then holds that state for 30~s while the runtime
watchdog polls every subsystem, so the calibration sweep already observes a loaded cell.

During training the active intent alternates every 600 control steps, i.e.\ every 60~s of
\emph{counted} control, giving 60 blocks and therefore 30 blocks (18{,}000 steps) per
direction. Between two blocks the controller does not pause: the runner spawns a worker that
rewrites the intent and waits for the A1 policy and rApp guidance to be regenerated, while
\texttt{run\_controller} keeps issuing actions under the phase label \texttt{intent\_switch}
with \texttt{max\_steps=None}. Those transitions are logged, and during training they are also
inserted into the replay, but they are excluded from the validated per-phase budget. The
wall-clock spacing between intent switches therefore exceeds 60~s. Empirically the surplus is
substantial: the asynchronous replay held 45{,}726 rows for the guided arm and 46{,}317 for
the unguided arm against 1{,}000 prefilled plus 36{,}000 training transitions, i.e.\
approximately $8.7\times10^3$ and $9.3\times10^3$ uncounted switch transitions per run.

Exploration noise is scheduled against the specification rather than the static
configuration: \texttt{exploration\_noise\_decay\_steps} is overridden to
$\operatorname{round}(0.4\times36{,}000)=14{,}400$, so Gaussian action noise decays linearly
from $0.2$ to $0.02$ over the first 40\% of training and leaves a 21{,}600-step low-noise
exploitation tail before evaluation.

Both evaluation phases are frozen. The runner sets the runtime state to
\texttt{train=False, publish=False}, then calls \texttt{freeze()} on the Actor-snapshot
watcher, which pins the serving Actor and sets \texttt{stats\_frozen}, so observation
normalization and reward scaling also stop updating. The freeze is verifiable in the artifact:
every arm reports exactly one distinct serving Actor version in each evaluation phase.

\subsection{Calibration}
Calibration is a fixed open-loop sweep, not a search and not a fourth policy. Slice~A is held
at each of 25\%, 35\%, 50\%, 65\%, and 75\% for 200 control steps (20~s) each, driven through
\texttt{experiment.fixed\_slice\_a\_ratio}, which makes the controller behave as a static
allocator regardless of arm. Three quantities are then derived.

\begin{enumerate}
  \item \emph{Cell capacity.} The runner sums \texttt{dl\_th\_mbps} over slices per
  \texttt{ts\_ms} in \texttt{network\_state} for the whole 100~s sweep and takes the
  linearly interpolated 95th percentile of those per-summary cell totals. The result is stored
  as \texttt{experiment.capacity\_mbps} and is the denominator of the normalized total and
  priority throughput terms in the reward, so every reward value is relative to this per-seed
  capacity estimate.
  \item \emph{Protected-slice floor.} Per-slice throughput samples are bucketed by the sweep
  level nearest to the acknowledged Slice-A ratio. For each intent the runner takes the
  reference allocation (65\% Slice-A for Intent~1, 35\% for Intent~2), computes the
  nearest-rank 10th percentile of the \emph{protected} slice throughput at that allocation,
  and sets $F=0.90\times p_{10}$. The floor is therefore a 10\% margin below an already
  pessimistic observed level, is per seed and per testbed, and is rendered into the
  natural-language intent string with two decimals through the \texttt{\{floor\_mbps\}}
  placeholder. Calibration raises and fails the run if $F\le0$ or if the reference allocation
  is not itself feasible.
  \item \emph{Guidance bounds.} A sweep level is labelled safe post hoc if its protected-slice
  $p_{10}$ meets $F$ and it respects the intent's direction. Writing $\mathcal{S}$ for that
  set, Intent~1 publishes Slice~A in $[55\%,\max\mathcal{S}]$ and Intent~2 publishes Slice~A
  in $[\min\mathcal{S},45\%]$, with complementary bounds for Slice~B. These are exported as
  A1 \texttt{policy\_constraints}.
\end{enumerate}

The 10--90\% envelope is \emph{not} a calibration product. It is the specification constant
pair \texttt{operational\_min\_prb}${}=10$ / \texttt{operational\_max\_prb}${}=90$, copied into
\texttt{control.min\_prb} / \texttt{control.max\_prb} before any run and applied by
\texttt{\_global\_bounds} to every arm. For the 106-PRB cell it gives the integer interval
$b_A\in[11,95]$. Only \texttt{llm\_only} and \texttt{llm\_guided\_ddpg} intersect the
calibration bounds with this envelope; \texttt{ddpg\_only} sees the envelope alone. With the
55\% minimum priority ratio the guided Intent-1 interval starts at
$b_A^{\min}=\max(\lceil 55\cdot106/100\rceil,\,106-\lfloor 45\cdot106/100\rfloor)=59$ and the
guided Intent-2 interval ends at $b_A^{\max}=47$, independently of $\mathcal{S}$.

The calibration profile is written to
\texttt{<scenario>/seed-<N>-calibration-profile.json} and reused by the remaining arms of the
same (scenario, seed) cell when the profile version matches. All three arms of a seed thus
share one capacity estimate, one floor, and one bound set, which is what makes the paired
comparison meaningful. Finally the 1{,}000 calibration transitions are replayed into the
learner buffer of the two learning arms with synthetic guidance that alternates every 100
steps, and the agent's SLA history, violation counter, and guided weight are reset before and
after the prefill so that calibration cannot bias the safety state machine.

\subsection{Traffic generation}
Each UE is served by one persistent iperf3 UDP client started inside the external data-network
container and addressed to the UE PDU IP, and by one iperf3 server bound to that PDU IP inside
the UE network namespace. There is exactly one stream and one port per UE; there is no
separate burst stream. The client duration is the larger of 3600~s and the specification key
\texttt{traffic\_max\_duration\_s}, i.e.\ 86{,}400~s, so that traffic never expires mid-run.

For dynamic scenarios the client is offered that UE's \emph{phase-maximum} rate for the whole
run and the phase pattern is realized in the network rather than at the source. The controller
installs an HTB qdisc on the external data-network egress device (\texttt{root handle 1:},
\texttt{r2q 100}, \texttt{default 999}, root and default classes at \texttt{1000mbit}), one
class \texttt{1:10}$x$ per UE with a \texttt{pfifo limit 100} child and a \texttt{u32} filter
matching the UE PDU address as a \texttt{/32}. Every 5~s the scheduler rewrites all five
classes with a single batched \texttt{tc class change ... htb rate Xmbit ceil Xmbit}. The
balanced scenario has a single phase, so no shaper is installed and the clients simply run at
the constant 30~Mbps per-UE rate; \texttt{traffic\_summary.json} records this as
\texttt{traffic\_backend}${}=$\texttt{iperf3} rather than \texttt{tc\_htb}.

Shaper fidelity is audited per segment. At the end of each UE $\times$ 5~s segment the
controller reads that UE's HTB class byte counter, computes the actual egress rate over the
elapsed interval, and accepts the segment only if it lasted at least $0.8\times5=4$~s and
delivered between 80\% and 125\% of the requested rate. The run-level
\texttt{dynamic\_segment\_success\_rate} divides valid segments by
$\lfloor D/5\rfloor$ (plus one if the remainder is at least 4~s) times five UEs, and the run is
rejected below 0.95. In the balanced pilot the shaper is never installed, so the reported
segment success rate of 100\% is vacuous and must not be read as evidence of shaper fidelity.

Receiver-side goodput is measured from the UE TUN interface \texttt{rx\_bytes} delta over the
measurement window. The reported loss percentage is
$100\,(1-\text{received}/\text{offered})$ where ``offered'' is the scenario's nominal mean
per-UE rate; it is an offered-versus-delivered figure computed at the receiver, not the
iperf3 sender's own datagram-loss report. Traffic phase applications are logged to a
\texttt{traffic\_events} table with planned and applied timestamps, and the trace is hashed so
that different arms can be shown to have seen identical offered load.

\subsection{Validity gates}
Table~\ref{tab:proto-gates} lists every check applied by \texttt{\_validate\_run}. A raising
check marks the run \texttt{failed}; the five-UE coverage check is non-fatal but marks the run
\texttt{degraded}. \texttt{scripts/build\_results.py} discards anything that is not
\texttt{status=complete} and \texttt{data\_quality=primary} before it opens a database, so
neither class can reach a table.

\begin{table*}[t]
\centering
\footnotesize
\caption{Validity gates. ``Spec'' checks read the campaign specification; ``code'' checks are
hardcoded in \texttt{experiment\_runner.py}. Async-only checks apply to \texttt{ddpg\_only}
and \texttt{llm\_guided\_ddpg}.}
\label{tab:proto-gates}
\begin{tabular}{llll}
\toprule
Check & Threshold & Source & Effect\\
\midrule
Per-phase transition count & exactly 1{,}000 / 36{,}000 / 1{,}800 / 1{,}800 & spec durations & fail\\
RC apply success rate & $\geq0.99$ (\texttt{min\_apply\_success\_rate}) & spec & fail\\
Mapped UEs & $\geq1$ under \texttt{allow\_ue\_churn}, else exactly 5 & spec & fail\\
Transitions with effect coverage $\geq0.8$ & $\geq95\%$ & code & fail\\
RC control latency p99 & $<20$~ms & code & fail\\
Controller tick jitter p95 / p99 & $\leq120$ / $\leq150$~ms & spec & fail (async)\\
Fresh-state action interval p95 & $\leq180$~ms & spec & fail (async)\\
Controller stale-skip rate & $\leq0.5$ & spec & fail (async)\\
Controller tick telemetry present & non-empty & code & fail (async)\\
Traffic summary parseable & --- & code & fail\\
Valid iperf UEs & $\geq1$ (\texttt{min\_valid\_traffic\_ues}) & spec & fail\\
Traffic event success rate & $\geq0.99$ & spec & fail\\
Dynamic segment success rate & $\geq0.95$ & spec & fail\\
Dynamic-channel command success / coverage & 100\% & code & fail (if enabled)\\
Five-UE coverage & $\geq0.90$ (\texttt{primary\_five\_ue\_coverage}) & spec & degrade\\
\bottomrule
\end{tabular}
\end{table*}

Two specification keys are computed and archived but never enforced and must not be described
as gates: \texttt{max\_fresh\_action\_interval\_p99\_ms}${}=300$ is recorded in the validation
block without an assertion, and \texttt{max\_summary\_age\_ms}${}=3000$ is never read by the
runner. One gate is also weaker than it appears: the mapped-UE query counts distinct
\texttt{ue\_id} values over the whole database without a run or time filter, so when several
runs share one database it does not isolate the run under test. Five-UE coverage, which is
computed strictly inside the run's own transition window, is the check that actually
constrains per-run UE presence.

Per-transition trainability is decided independently of these run-level gates. A transition
enters the replay only if the control was acknowledged, at least one UE was active, the number
of mapped UEs equals the number of active UEs, no counter reset occurred in the window, all
metric sources were valid, and the action-effect window coverage was at least 0.8.

\subsection{Reported metrics and their definitions}
All phase metrics are computed by \texttt{analyze\_results.\_step\_metrics} over the
transitions of a phase. Two definitions carry most of the interpretation and are used
throughout Section~\ref{sec:results-anchor}:

\begin{itemize}
  \item \emph{SLA satisfaction} is the fraction of transitions whose normalized protected-slice
  deficit is zero, i.e.\ the protected slice met its floor $F$ (or the floor was inapplicable
  because $F\le0$ or the protected slice carried no UE).
  \item \emph{Intent satisfaction} is the fraction of transitions that satisfy the SLA
  \emph{and} command the priority slice at least its minimum ratio of 55\%. It is therefore a
  strictly stronger condition, and the gap between the two is a commanded-allocation property,
  not a throughput shortfall.
\end{itemize}

The remaining reported quantities are: SLA deficit area (the sum of positive normalized
deficits over a phase), protected-slice 5th-percentile throughput, aggregate and per-slice
RLC-derived throughput, UE-level and slice-level Jain fairness, commanded and KPM-measured PRB
share, RLC buffer occupancy, receiver goodput and offered-versus-delivered loss, RC control
latency, controller tick jitter, stale-state skip rate, learner TD error, and Actor
diagnostics. The Actor diagnostics are: saturation rate (fraction of raw actor outputs at or
beyond $0.02$/$0.98$), action-bin coverage (visited tenths of the unit action range, out of
ten), and acknowledged-action boundary rate (fraction of commanded Slice-A budgets within half
a PRB of an edge of that transition's own feasible interval). Early-phase diagnostics use the
first 300 training transitions; ``reward AUC'' is the sum of raw rewards over that window.
Downlink BLER and wideband CQI are recorded but are effectively unpopulated under the
\texttt{static\_awgn} RFSimulator profile used here, and are consequently not interpreted as
measured radio quality.

\ifFullCampaign
The complete campaign uses same-seed paired bootstrap confidence intervals.
\else
The complete campaign will use same-seed paired bootstrap confidence intervals.
\fi
A single-seed pilot cannot estimate run-to-run uncertainty; repeated windows inside one run
are not treated as independent seeds.

\subsection{Reproducibility}
\textbf{Seeding.} Each run calls \texttt{random.seed(seed)} in the runner, and the agent seeds
both Python's RNG and \texttt{torch} with the run seed in the controller process and again in
the learner subprocess, which is launched with \texttt{--seed}. Dynamic-channel traces, when
enabled, are drawn with derived seeds $s$, $s+10^4$, and $s+2\times10^4$ for the training and
two evaluation phases. Arm execution order inside a (seed, scenario) cell is shuffled
deterministically from the string \texttt{"seed:scenario"} so that arm order is not confounded
with time. NumPy, CUDA kernels, and cuDNN determinism flags are not seeded or forced, and the
learner advances on wall-clock rather than in lockstep with the controller; runs are therefore
statistically, not bitwise, reproducible.

\textbf{Isolation.} Every run gets its own result directory containing its checkpoint
(\texttt{ddpg.pt}), Actor-snapshot directory, effective configuration, learner log, archived
component logs, traffic summary and shaper-segment file, manifest, and an end-of-run SQLite
backup of the database. The live database is \emph{not} per run: it is the single shared file
\texttt{/tmp/llm\_hric/llm\_hric.sqlite3}. It is deleted, together with its WAL/SHM files and
the actuator socket, only inside \texttt{restart\_stack}, which the runner invokes only when
started with \texttt{--manage-stack}; that flag also tears down and relaunches the whole
RFSimulator stack (\texttt{UE\_MODE=multi}, \texttt{UE\_COUNT=5}, monitor and guidance
services enabled, standalone DDPG disabled) before every arm/seed and then waits up to 180~s
for preflight. Campaign runs must therefore be launched with \texttt{--manage-stack};
without it, a run appends to whatever database is already present.

\textbf{Resume and failure handling.} \texttt{--resume} scans archived manifests and skips any
(scenario, arm, seed) tuple already recorded as complete or degraded, making the campaign
idempotent across interruptions; frozen-policy re-evaluation manifests are ignored by this
scan, and dynamic-scenario runs produced by the superseded host-side termination path are
deliberately not counted as complete and are retried. \texttt{--fail-fast} aborts after
archiving the first failed run; otherwise failures are accumulated and re-raised at the end so
that one bad cell does not cost the rest of the campaign. Preflight is mandatory and requires
a GPU unless \texttt{--allow-no-gpu} is given.

\textbf{Cost.} Measured from the run manifests, a single run takes 5{,}113--5{,}276~s
(85.2--87.9~min) of wall clock against a 4{,}090~s nominal phase budget. The roughly 27\%
overhead is dominated by the 64 intent switches per run (two after calibration, 60 between
training blocks, and one before each evaluation), each of which waits for a fresh rApp
inference at the 10~s rApp period, plus stack settling and traffic prefill. At the observed
mean of 5{,}210~s per run, the complete $3\times3\times5$ campaign requires approximately
65~hours of continuous testbed time excluding stack restarts.

\textbf{Environment.} Table~\ref{tab:proto-env} lists the testbed. The single 8~GB laptop GPU
hosts both the 4-bit quantized Gemma rApp and the TD3-style networks; the controller is
restricted to one Torch thread and the learner to two so that inference cannot starve the
100~ms control loop. The runner records the GPU model, driver, Python, and Torch/Transformers
versions into every manifest, so a single future run restores this block exactly; the values
below were probed on the testbed host because the pilot manifests were stored under
\texttt{/tmp} and are no longer available.

\begin{table}[t]
\centering
\footnotesize
\caption{Testbed hardware and software.}
\label{tab:proto-env}
\begin{tabular}{ll}
\toprule
Item & Value\\
\midrule
GPU & NVIDIA GeForce RTX 4060 Laptop, 8188~MiB, driver 580.173.02\\
CPU & Intel Core i7-13650HX, 1 socket, 14 cores, 20 threads\\
RAM & 31~GiB\\
OS & Ubuntu 22.04.5 LTS, kernel 6.5.0-18-generic\\
PyTorch & 2.7.0+cu118\\
OAI revision & \texttt{cb0e5012}\\
FlexRIC revision & \texttt{340c36bc}\\
\bottomrule
\end{tabular}
\end{table}

\textbf{Artifact custody.} The primary evidence for every number in
Section~\ref{sec:results-anchor} is the per-run SQLite backup (1.0--1.3~GB per run) plus its
manifest. These files were written under \texttt{/tmp} in the pilot and have since been lost
with the host's temporary directory; the committed \texttt{generated/} CSV and \LaTeX{}
artifacts are what remain. Future campaign runs must therefore use a \texttt{--results}
directory outside \texttt{/tmp}, and the databases must be archived with the paper.
